% ============================================================
%  CHAPITRE II : Cadre théorique et état de l'art
% ============================================================
\chapter{Cadre théorique et état de l'art}
\label{chap:etat_art}

\section*{Introduction}

Ce chapitre rassemble les fondements théoriques mobilisés dans le projet. Il
rappelle d'abord le rôle et le coût des stocks, puis présente les méthodes de
segmentation du portefeuille de références, les modèles classiques de
réapprovisionnement, le dimensionnement du stock de sécurité, les techniques de
prévision de la demande et enfin les indicateurs de performance retenus.

% -------------------------------------------------------
\section{Rôle et coût des stocks}
\label{sec:role_stocks}

\subsection{Fonctions du stock}

\iacsdefinition{Stock}{
  Ensemble des biens détenus par l'entreprise en attente d'utilisation ou de
  vente. Le stock joue un rôle de tampon entre des flux d'approvisionnement et
  de consommation qui ne sont ni synchronisés ni parfaitement prévisibles.
}

On distingue classiquement plusieurs composantes :

\begin{itemize}
  \item le \textbf{stock de cycle}, qui résulte du fait que les
    approvisionnements se font par lots plutôt qu'à l'unité ;
  \item le \textbf{stock de sécurité}, destiné à absorber les aléas de la
    demande et des délais de livraison ;
  \item le \textbf{stock en transit}, déjà commandé mais non encore réceptionné ;
  \item le \textbf{stock mort ou dormant}, sans consommation sur une période
    prolongée.
\end{itemize}

\subsection{Structure des coûts}

Le tableau~\ref{tab:couts_stock} présente les trois familles de coûts qui
s'opposent dans tout arbitrage de gestion de stock.

\begin{table}[H]
  \centering
  \caption{Familles de coûts liées à la gestion des stocks}
  \label{tab:couts_stock}
  \renewcommand{\arraystretch}{1.4}
  \begin{tabularx}{\textwidth}{|l|X|c|}
    \hline
    \rowcolor{iacsTableHead}
    \textcolor{white}{\textbf{Famille}} &
    \textcolor{white}{\textbf{Contenu}} &
    \textcolor{white}{\textbf{Notation}} \\
    \hline
    Coût de passation &
    Traitement de la commande, transport, réception, contrôle &
    $S$ \\
    \hline
    \rowcolor{iacsLightGray}
    Coût de possession &
    Immobilisation financière, stockage, assurance, obsolescence &
    $H$ \\
    \hline
    Coût de rupture &
    Arrêt d'exploitation, achat en urgence, dégradation du service &
    $C_{r}$ \\
    \hline
  \end{tabularx}
\end{table}

Le coût de possession est généralement exprimé en pourcentage de la valeur
unitaire de l'article : $H = i \cdot c$, où $i$ est le taux de possession annuel
et $c$ le coût unitaire.

\iacswarning{Le coût de rupture est le plus difficile à estimer et le plus
souvent sous-évalué. Dans un contexte industriel, il correspond au coût horaire
d'immobilisation de l'installation, sans commune mesure avec la valeur de
l'article manquant.}

% -------------------------------------------------------
\section{Segmentation du portefeuille de références}
\label{sec:segmentation}

\subsection{Classification ABC}

La classification \gls{abc} repose sur le principe de Pareto : une faible part
des références représente l'essentiel de la valeur consommée. Les références
sont triées par valeur de consommation annuelle décroissante, puis réparties en
trois classes selon la valeur cumulée.

\begin{equation}
  V_{i} = D_{i} \cdot c_{i}
  \qquad\text{et}\qquad
  P_{k} = \frac{\sum_{i=1}^{k} V_{(i)}}{\sum_{i=1}^{N} V_{(i)}}
  \label{eq:abc}
\end{equation}

où $V_{i}$ est la valeur de consommation annuelle de la référence $i$, $D_{i}$
sa consommation annuelle, $c_{i}$ son coût unitaire, et $P_{k}$ la part cumulée
après tri décroissant.

\begin{table}[H]
  \centering
  \caption{Seuils usuels de la classification ABC}
  \label{tab:seuils_abc}
  \renewcommand{\arraystretch}{1.4}
  \begin{tabularx}{\textwidth}{|c|c|c|X|}
    \hline
    \rowcolor{iacsTableHead}
    \textcolor{white}{\textbf{Classe}} &
    \textcolor{white}{\textbf{\% valeur}} &
    \textcolor{white}{\textbf{\% références}} &
    \textcolor{white}{\textbf{Mode de gestion}} \\
    \hline
    A & $\approx 80\%$ & $\approx 20\%$ & Suivi rapproché, taux de service élevé \\
    \hline
    \rowcolor{iacsLightGray}
    B & $\approx 15\%$ & $\approx 30\%$ & Gestion sur règles automatiques \\
    \hline
    C & $\approx 5\%$  & $\approx 50\%$ & Gestion allégée, lots de réapprovisionnement larges \\
    \hline
  \end{tabularx}
\end{table}

\subsection{Classification XYZ}

La classification \gls{xyz} complète la précédente en caractérisant la
\emph{régularité} de la demande plutôt que sa valeur. Elle s'appuie sur le
coefficient de variation :

\begin{equation}
  CV = \frac{\sigma_{d}}{\bar{d}}
  \label{eq:cv}
\end{equation}

où $\bar{d}$ est la demande moyenne par période et $\sigma_{d}$ son écart-type.
Les références sont classées X (demande régulière, $CV$ faible), Y (demande
variable ou saisonnière) ou Z (demande erratique, $CV$ élevé).

Le croisement des deux classifications produit une matrice à neuf cases qui
oriente le mode de gestion : les références AX se prêtent à un
réapprovisionnement automatique à stock de sécurité réduit, tandis que les
références AZ, coûteuses et imprévisibles, exigent un suivi manuel.

\iacsnote{Le croisement ABC/XYZ constitue le point d'entrée de la démarche : il
évite d'appliquer une règle unique à un portefeuille hétérogène.}

% -------------------------------------------------------
\section{Modèles de réapprovisionnement}
\label{sec:modeles_appro}

\subsection{Quantité économique de commande}

Le modèle de Wilson détermine la quantité de commande $Q$ qui minimise la somme
du coût de passation et du coût de possession, sous hypothèse de demande
constante et de délai fixe. Le coût total annuel s'écrit :

\begin{equation}
  C(Q) = \frac{D}{Q} \cdot S + \frac{Q}{2} \cdot H
  \label{eq:cout_total}
\end{equation}

L'annulation de la dérivée $\dfrac{dC}{dQ} = -\dfrac{DS}{Q^{2}} + \dfrac{H}{2}$
conduit à la \gls{eoq} :

\begin{equation}
  Q^{*} = \sqrt{\frac{2DS}{H}}
  \label{eq:eoq}
\end{equation}

\iacswarning{Les hypothèses du modèle de Wilson (demande déterministe et
constante, délai fixe, absence de remise sur quantité) sont rarement vérifiées
en pratique. Le résultat reste néanmoins un ordre de grandeur utile, la fonction
de coût étant peu sensible autour de son optimum.}

\subsection{Politiques de gestion}

Deux politiques principales sont utilisées :

\begin{itemize}
  \item la politique \textbf{$(s, Q)$} — à point de commande : une quantité
    fixe $Q$ est commandée dès que le stock disponible franchit le seuil $s$ ;
  \item la politique \textbf{$(T, S)$} — à recomplètement périodique : le stock
    est ramené à un niveau cible $S$ à intervalles réguliers $T$.
\end{itemize}

\subsection{Point de commande}

Le point de commande \gls{rop} correspond au niveau de stock déclenchant le
réapprovisionnement. Il couvre la consommation attendue pendant le délai de
livraison, augmentée du stock de sécurité :

\begin{equation}
  ROP = \bar{d} \cdot L + SS
  \label{eq:rop}
\end{equation}

où $\bar{d}$ est la demande moyenne par période et $L$ le délai de
réapprovisionnement exprimé dans la même unité.

% -------------------------------------------------------
\section{Dimensionnement du stock de sécurité}
\label{sec:stock_securite}

\subsection{Taux de service et coefficient de sécurité}

\iacsdefinition{Taux de service (cycle service level)}{
  Probabilité de ne pas connaître de rupture au cours d'un cycle de
  réapprovisionnement. À un taux de service cible $\alpha$ correspond un
  coefficient $z = \Phi^{-1}(\alpha)$, quantile de la loi normale centrée
  réduite.
}

\begin{table}[H]
  \centering
  \caption{Coefficient de sécurité en fonction du taux de service}
  \label{tab:coef_z}
  \renewcommand{\arraystretch}{1.4}
  \begin{tabular}{|c|c|c|c|c|c|c|}
    \hline
    \rowcolor{iacsTableHead}
    \textcolor{white}{\textbf{$\alpha$}} &
    \textcolor{white}{90\%} & \textcolor{white}{95\%} &
    \textcolor{white}{97,5\%} & \textcolor{white}{98\%} &
    \textcolor{white}{99\%} & \textcolor{white}{99,9\%} \\
    \hline
    $z$ & 1,28 & 1,65 & 1,96 & 2,05 & 2,33 & 3,09 \\
    \hline
  \end{tabular}
\end{table}

\subsection{Formules de calcul}

Lorsque seule la demande est incertaine et que le délai est constant :

\begin{equation}
  SS = z \cdot \sigma_{d} \cdot \sqrt{L}
  \label{eq:ss_demande}
\end{equation}

Lorsque la demande et le délai sont tous deux incertains, la formule combinée
s'écrit :

\begin{equation}
  SS = z \cdot \sqrt{ L \cdot \sigma_{d}^{2} + \bar{d}^{\,2} \cdot \sigma_{L}^{2} }
  \label{eq:ss_combine}
\end{equation}

où $\sigma_{d}$ est l'écart-type de la demande par période, $\sigma_{L}$ celui
du délai de livraison, $\bar{d}$ la demande moyenne et $L$ le délai moyen.

\iacsnote{L'équation~\eqref{eq:ss_combine} montre que le stock de sécurité croît
comme la racine carrée du délai : \emph{réduire et fiabiliser les délais
fournisseurs est souvent plus efficace que d'augmenter le stock}.}

\subsection{Cas des demandes intermittentes}

Pour les références à demande rare — fréquentes parmi les pièces de rechange —
l'hypothèse de normalité n'est plus valable. La demande est alors mieux décrite
par une loi de Poisson ou par les méthodes dédiées à la demande intermittente
(méthode de Croston et ses variantes).

% -------------------------------------------------------
\section{Prévision de la demande}
\label{sec:prevision}

\subsection{Méthodes retenues}

\begin{table}[H]
  \centering
  \caption{Méthodes de prévision envisagées}
  \label{tab:methodes_prevision}
  \renewcommand{\arraystretch}{1.4}
  \begin{tabularx}{\textwidth}{|l|X|l|}
    \hline
    \rowcolor{iacsTableHead}
    \textcolor{white}{\textbf{Méthode}} &
    \textcolor{white}{\textbf{Principe}} &
    \textcolor{white}{\textbf{Usage}} \\
    \hline
    Moyenne mobile &
    Moyenne des $k$ dernières périodes &
    Référence de base \\
    \hline
    \rowcolor{iacsLightGray}
    Lissage exponentiel (\gls{ses}) &
    Pondération décroissante du passé, paramètre $\alpha$ &
    Demande stable \\
    \hline
    Holt-Winters &
    Lissage avec tendance et saisonnalité &
    Demande saisonnière \\
    \hline
    \rowcolor{iacsLightGray}
    Croston &
    Séparation intervalle / taille de demande &
    Demande intermittente \\
    \hline
  \end{tabularx}
\end{table}

Le lissage exponentiel simple s'écrit :

\begin{equation}
  \hat{d}_{t+1} = \alpha \, d_{t} + (1 - \alpha) \, \hat{d}_{t},
  \qquad 0 < \alpha < 1
  \label{eq:ses}
\end{equation}

\subsection{Mesure de la qualité des prévisions}

\begin{equation}
  MAE = \frac{1}{n}\sum_{t=1}^{n} \left| d_{t} - \hat{d}_{t} \right|
  \qquad
  RMSE = \sqrt{\frac{1}{n}\sum_{t=1}^{n} \left( d_{t} - \hat{d}_{t} \right)^{2}}
  \label{eq:erreurs}
\end{equation}

\begin{equation}
  MAPE = \frac{100}{n}\sum_{t=1}^{n}
    \left| \frac{d_{t} - \hat{d}_{t}}{d_{t}} \right|
  \label{eq:mape}
\end{equation}

\iacswarning{Le \gls{mape} n'est pas défini lorsque $d_{t} = 0$, situation
fréquente pour les références à demande intermittente. Le \gls{mae} lui est
préféré dans ce cas.}

% -------------------------------------------------------
\section{Indicateurs de performance}
\label{sec:kpi}

\begin{table}[H]
  \centering
  \caption{Indicateurs de performance retenus}
  \label{tab:kpis}
  \renewcommand{\arraystretch}{1.5}
  \begin{tabularx}{\textwidth}{|l|c|X|}
    \hline
    \rowcolor{iacsTableHead}
    \textcolor{white}{\textbf{Indicateur}} &
    \textcolor{white}{\textbf{Formule}} &
    \textcolor{white}{\textbf{Interprétation}} \\
    \hline
    Taux de rotation &
    $R = \dfrac{\text{Consommation}}{\text{Stock moyen}}$ &
    Nombre de renouvellements du stock par an \\
    \hline
    \rowcolor{iacsLightGray}
    Couverture &
    $C = \dfrac{365}{R}$ &
    Nombre de jours de consommation couverts \\
    \hline
    Taux de service &
    $\dfrac{\text{Lignes servies}}{\text{Lignes demandées}}$ &
    Capacité à répondre à la demande \\
    \hline
    \rowcolor{iacsLightGray}
    Taux de rupture &
    $1 - \text{taux de service}$ &
    Part des demandes non satisfaites \\
    \hline
    Valeur du stock &
    $\sum_{i} q_{i} \cdot c_{i}$ &
    Capital immobilisé \\
    \hline
    \rowcolor{iacsLightGray}
    Stock dormant &
    Part sans mouvement sur 12 mois &
    Gisement d'économies immédiat \\
    \hline
  \end{tabularx}
\end{table}

% -------------------------------------------------------
\section{Travaux similaires}
\label{sec:travaux_similaires}

\hspace{1cm}[Présenter deux à quatre travaux comparables — articles, mémoires ou
retours d'expérience industriels portant sur l'optimisation des stocks de pièces
de rechange — et positionner la présente démarche par rapport à eux : périmètre,
méthodes employées, résultats annoncés.]

\section*{Conclusion}

Ce chapitre a rassemblé les outils théoriques nécessaires : segmentation
\gls{abc}/\gls{xyz}, modèles de réapprovisionnement, dimensionnement du stock de
sécurité et indicateurs de suivi. Le chapitre suivant décrit leur mise en
\oe uvre sur les données réelles de l'\gls{ocp}.
